\title{Emerging Trends in AI Agent Evaluation}
\begin{document}
\maketitle

\begin{abstract}
AI models are increasingly deployed as agents: rather than responding to a single prompt, they plan over many steps, call external tools, and act within software and real-world environments. This shift challenges the existing outcome-focused benchmarking paradigm for evaluating AI models. Evaluating these agent systems must now encompass an array of scaffolds, tools, and environments that present a far more complex challenge for evaluators today. First, this shift has necessitated many improvements to measurement methodology, scientific rigor, and evaluation infrastructure. Second, the deployment of agents in real settings has expanded the scope of evaluation from the lab to the field. Third, the accompanying changes within the broader evaluator ecosystem hold new implications for those interested in contributing to this work. We survey 33 shifts, practices, and questions driven by the AI agent paradigm across the above three dimensions—the “How,” “What,” and “Who” of AI agent evaluation. Our goal is to help consolidate the emerging discipline of agent evaluation, equipped with an evolving methodological toolkit (Part 1), an orienting research agenda (Part 2), and an updated portrait of the evaluation landscape (Part 3). Part 1: How Evaluation science and methods 1a Validity 1b Benchmark quality 1c Metrics beyond accuracy 1d Statistical rigor 1e Causal attribution 1f Contamination 1g Dynamic benchmarks 1h LLM judges and monitors 1i Log analysis 1j Evaluation infrastructure 1k
\end{abstract}

\section{Recovered Text}
Emerging Trends in AI Agent Evaluation
Felix Chen1,*
Stephan Rabanser1,*
Peter Kirgis1
Nitya Nadgir1
Rishi Bommasani2
Magda Dubois3
Elham Tabassi4
Cozmin Ududec3
Anka Reuel2
Sanmi Koyejo2
Sayash Kapoor1,5,*
Arvind Narayanan1,6,*
1Princeton University
2Stanford University
3UK AI Security Institute
4The Brookings Institution
5UC Berkeley
6Institute for Advanced Study
*Corresponding authors: \{felix\_chen, rabanser, arvindn\}@princeton.edu, sayashk@berkeley.edu
October 19, 2026
Abstract
AI models are increasingly deployed as agents: rather than responding to a single prompt, they plan
over many steps, call external tools, and act within software and real-world environments. This shift
challenges the existing outcome-focused benchmarking paradigm for evaluating AI models. Evaluating
these agent systems must now encompass an array of scaffolds, tools, and environments that present a
far more complex challenge for evaluators today. First, this shift has necessitated many improvements
to measurement methodology, scientific rigor, and evaluation infrastructure. Second, the deployment
of agents in real settings has expanded the scope of evaluation from the lab to the field. Third, the
accompanying changes within the broader evaluator ecosystem hold new implications for those interested
in contributing to this work. We survey 33 shifts, practices, and questions driven by the AI agent paradigm
across the above three dimensions—the “How,” “What,” and “Who” of AI agent evaluation. Our goal is to
help consolidate the emerging discipline of agent evaluation, equipped with an evolving methodological
toolkit (Part 1), an orienting research agenda (Part 2), and an updated portrait of the evaluation
landscape (Part 3).
Part 1: How
Evaluation science and methods
1a
Validity
1b
Benchmark quality
1c
Metrics beyond accuracy
1d
Statistical rigor
1e
Causal attribution
1f
Contamination
1g
Dynamic benchmarks
1h
LLM judges and monitors
1i
Log analysis
1j
Evaluation infrastructure
1k
Red-teaming
1l
Open-world evaluations
1m Human baselines
each trend closes with
best practices
Part 2: What
Expanded research agenda
2a Simulated users
2b Lab experiments
2c Field experiments
2d Observational usage data
2e Domain-specific evaluation
2f
Economic impact
2g Behavior under evaluation
2h Dangerous capabilities
2i
Security
2j
Interaction harms
2k Systemic evaluation
each trend closes with
research directions
Part 3: Who
Evaluator ecosystem
3a Who evaluates
3b Independence
3c Labs shaping evaluations
3d Evaluations shaping labs
3e Funding
3f
Rising costs
3g Venues
3h Public sector
3i
Disciplinary entry points
each trend closes with
practical considerations
Figure 1: Overview of the 33 trends surveyed in this paper.
1
Contents
Introduction
3
Scope and methods
6
Intended audience . . . . . . . . . . . . . . . . . . . . . . . . . . . . . . . . . . . . . . . . . . . . . . . . .
6
Agents as the object of evaluation . . . . . . . . . . . . . . . . . . . . . . . . . . . . . . . . . . . . . . . . .
6
Methods . . . . . . . . . . . . . . . . . . . . . . . . . . . . . . . . . . . . . . . . . . . . . . . . . . . . . . .
8
Part 1: Improved Evaluation Science and Methods (“How”)
8
1a. How does validity change for agents? . . . . . . . . . . . . . . . . . . . . . . . . . . . . . . . . . . . . .
8
1b. How should we evaluate benchmark quality?
. . . . . . . . . . . . . . . . . . . . . . . . . . . . . . . . 10
1c. What should agent evaluations measure beyond accuracy? . . . . . . . . . . . . . . . . . . . . . . . . . 11
1d. How can agent evaluations achieve statistical rigor? . . . . . . . . . . . . . . . . . . . . . . . . . . . . . 12
1e. How can agent evaluations attribute causal changes in performance? . . . . . . . . . . . . . . . . . . . 14
1f. How can agent evaluations resist contamination?
. . . . . . . . . . . . . . . . . . . . . . . . . . . . . . 15
1g. What role do dynamic benchmarks play, and what are their tradeoffs?
. . . . . . . . . . . . . . . . . . 16
1h. Where should LLMs judge and monitor agents? . . . . . . . . . . . . . . . . . . . . . . . . . . . . . . . 17
1i. What can log analysis reveal beyond outcomes? . . . . . . . . . . . . . . . . . . . . . . . . . . . . . . . 18
1j. What is the value of standardized evaluation infrastructure for agents?
. . . . . . . . . . . . . . . . . . 19
1k. How should red-teaming evolve for agents? . . . . . . . . . . . . . . . . . . . . . . . . . . . . . . . . . 20
1l. What can open-world evaluations reveal that benchmarks cannot? . . . . . . . . . . . . . . . . . . . . . 21
1m. How can evaluators interpret agent performance relative to human capabilities?
. . . . . . . . . . . . 21
Part 2: Expanding the Research Agenda for Evaluation (“What”)
23
2a. Human-agent interaction: evaluation with simulated users . . . . . . . . . . . . . . . . . . . . . . . . . 24
2b. Human-agent interaction: experiments in the lab . . . . . . . . . . . . . . . . . . . . . . . . . . . . . . 25
2c. Human-agent interaction: field experiments . . . . . . . . . . . . . . . . . . . . . . . . . . . . . . . . . 26
2d. Observational data on real-world usage
. . . . . . . . . . . . . . . . . . . . . . . . . . . . . . . . . . . 27
2e. Domain-specific evaluation
. . . . . . . . . . . . . . . . . . . . . . . . . . . . . . . . . . . . . . . . . . 27
2f. Economic impact evaluation . . . . . . . . . . . . . . . . . . . . . . . . . . . . . . . . . . . . . . . . . . 28
2g. Agent behavior under evaluation . . . . . . . . . . . . . . . . . . . . . . . . . . . . . . . . . . . . . . . 30
2h. Dangerous capabilities and risk . . . . . . . . . . . . . . . . . . . . . . . . . . . . . . . . . . . . . . . . 31
2i. Agent and evaluation security . . . . . . . . . . . . . . . . . . . . . . . . . . . . . . . . . . . . . . . . . 32
2j. Interaction harms . . . . . . . . . . . . . . . . . . . . . . . . . . . . . . . . . . . . . . . . . . . . . . . . 32
2k. Systemic evaluation
. . . . . . . . . . . . . . . . . . . . . . . . . . . . . . . . . . . . . . . . . . . . . . 34
Part 3: The Evolving Evaluator Ecosystem (“Who”)
34
3a. Who evaluates? . . . . . . . . . . . . . . . . . . . . . . . . . . . . . . . . . . . . . . . . . . . . . . . . . 35
3b. What does independence require?
. . . . . . . . . . . . . . . . . . . . . . . . . . . . . . . . . . . . . . 36
3c. Does lab direction drive evaluations? . . . . . . . . . . . . . . . . . . . . . . . . . . . . . . . . . . . . . 37
3d. Do evaluations drive lab direction? . . . . . . . . . . . . . . . . . . . . . . . . . . . . . . . . . . . . . . 38
3e. How does funding shape independent evaluation?
. . . . . . . . . . . . . . . . . . . . . . . . . . . . . 39
3f. How do rising costs affect evaluations? . . . . . . . . . . . . . . . . . . . . . . . . . . . . . . . . . . . . 40
3g. Where are the venues for evaluations? . . . . . . . . . . . . . . . . . . . . . . . . . . . . . . . . . . . . 41
3h. How is the public sector responding? . . . . . . . . . . . . . . . . . . . . . . . . . . . . . . . . . . . . . 42
3i. Where are the entry points for different disciplines? . . . . . . . . . . . . . . . . . . . . . . . . . . . . . 43
Conclusion
45
Acknowledgements
45
A
Best Practices, Research Directions, and Practical Considerations
66
2
Introduction
For much of the history of artificial intelligence (AI) and machine learning (ML) models, evaluation followed
a simple template: query a model, collect its response, and compare that response against a given reference
answer. However, with the advent of AI agents, evaluation has grown into a significantly more complex
enterprise (see Figure 2). The basic mental model of evaluation as outlined above belies the sprawling set
of resources, infrastructure, and design considerations required by evaluators today. Questions of interest
for evaluation of AI agents today no longer implicate models alone, but rather expand to the larger systems
around them, consisting of scaffolds (also referred to as “agent harnesses”) [16, 116, 138], tools, testing
environments, and user interfaces. As a result, long-standing concerns may manifest differently for agents
than they did for models, and previous assumptions may not hold (see Table 1). These changes in the
overall evaluation landscape are taking place amidst rising expectations for AI agents and demands for
longer-horizon tasks. Meeting these demands requires updating our understanding of principled AI agent
evaluation in three key dimensions (see Figure 1): how to evaluate, what to evaluate, and who evaluates.
1) Rethink evaluation methodology (“How?”). First, agent evaluation requires revisiting and imagining
evaluation methods. Conceptually, familiar concerns about evaluation take new form within the agentic
paradigm. For example, establishing validity becomes harder when performance depends not only on the
underlying model, but also on the particular agent scaffold and environment. Achieving statistical rigor must
better account for long, stochastic trajectories, where small differences early on can compound into sharply
divergent outcomes. Concerns about contamination extend beyond training data alone; agents now have
access to a wide array of tools and web access that make identifying such instances harder. Methodologically,
these shifts create new opportunities for metrics of interest and evaluation infrastructure. Conducting
evaluations today requires an elaborate set of harnesses, autograders, and sandboxed environments, which
are themselves new surfaces for scrutiny. We offer a set of emerging best practices.
2) Broaden the scope of evaluation (“What?”). Second, the agentic shift demands a broadened research
agenda for evaluators. Agent deployment in real settings has expanded the scope of evaluation from the
lab to the field, necessitating coverage of capabilities, human interactions, and systemic considerations.
No longer can questions for evaluation be confined to a question of capabilities in a vacuum, but must
also examine how agents are situated in new relationships with humans and each other. Fueled by public
debates around the economic and societal impact of AI, we separate evaluations that ask whether an
agent can do a task on its own (automation) from those that ask whether they can support human work
(augmentation). We offer a set of relevant research directions.
3) Recognize the role of evaluators (“Who?”). Third, evaluators shape the evidence that developers,
policymakers, and the general public rely on. Evaluators influence the broader discourse not only by setting
questions of interest, but by determining how evidence is produced, analyzed, and shared. Who designs
an evaluation therefore matters because the choices evaluators make about scaffolds, prompts, tool use,
environments, budgets, points of human intervention, and statistical analyses all create substantial degrees
of freedom. Consequently, the historic understanding of benchmarks as a source of evaluator-agnostic
results can no longer be taken for granted. Reflecting the changing landscape of this work, we offer practical
considerations for researchers interested in this work.
There is an urgent need to identify and disseminate best practices, chart a path forward for the agent
research community, and develop an ecosystem capable of conducting these evaluations. This paper is a
step in that direction. We begin by clarifying its scope and methodology. Part 1 examines the conceptual
concerns, emerging tools, and evolving protocols that are shaping a science of AI agent evaluation. Part 2
turns to the widened scope of evaluation, emphasizing both technical and sociotechnical lenses, micro-
and macro-levels of analysis, and evolving conceptions of safety, security, and harm. Part 3 characterizes
3
the current evaluation ecosystem, and discusses how changes in evaluators’ material conditions and
independence impact the choices available.
Table 1: Why agent evaluation is a new field and why LLM evaluation practices are insufficient. Practical suggestions
for the evaluator ecosystem are treated separately.
Methodological Question or Area of
Evaluation
How Agents Exacerbate Existing Model Evaluation Challenges, or
Create New Ones
1a. How does validity change for
agents?
The outcome-oriented focus of model evaluations breaks down with
agents, since they can reach the same outcomes through disparate
trajectories, possibly via shortcuts. Moreover, constructs of interest are
harder to operationalize due to the dynamic nature of agents’ real-world
interactions. Evaluators must be more precise about constructs, metrics,
and how they think about validity.
1b. How should we evaluate
benchmark quality?
Manual human review of model benchmark questions was more feasible,
and impossible tasks more clear to verify—with the increase in
components for agent evaluation tasks, it is far more tedious to attempt
validating broken or impossible agent tasks.
1c. What should agent evaluations
measure beyond accuracy?
An outcome-oriented evaluation mindset for models prioritized accuracy
(i.e., success on tasks) as the key metric of interest. But agent
performance includes far more than just accuracy; deployment in the real
world involves tradeoffs between cost, latency, and reliability, important
concerns not captured by accuracy alone.
1d. How can agent evaluations
achieve statistical rigor?
Agents create greater variance across runs and over multi-step, possibly
disparate, trajectories; this increased variance demands better ways to
capture and model this uncertainty.
1e. How can agent evaluations
attribute causal changes in
performance?
Model evaluation benefited from clear causal attribution—changes in
models drove changes in outputs. With a larger number of sources of
variance (e.g., environments, tools, scaffolds, etc) agents now have more
possible potential causal sources to rule out when agent performance
changes. Moreover, the same scaffold may not interact with all models in
the same way, complicating the separation of model from scaffold
attribution.
1f. How can agent evaluations resist
contamination?
Contamination for model evaluation largely was conceptualized in terms
of prior exposure to benchmark content. With agents and their access to
scaffolds and tool use, there are now far more possible sources of
contamination.
1g. What role do dynamic benchmarks
play, and what are their tradeoffs?
Dynamic benchmarks offer one response to contamination and saturation,
but their distribution shifts are known to create comparability and
coverage issues. This dynamic nature creates another variance source for
agent evaluations, making attributions of score changes even harder.
1h. Where should LLMs judge and
monitor agents?
Agent evaluations have dramatically amplified the need and utility of
LLM judges and monitoring thanks to long trajectories and open-ended
outputs that are outpacing manual human review. The increasing
centrality of these setups is creating fresh risks around overreliance and
questions about their efficacy and validity. Moreover, it is raising questions
of when—and where—we need humans in the loop to judge and monitor
agents? What information do LLMs or humans need for that task?
4
Methodological Question or Area of
Evaluation
How Agents Exacerbate Existing Model Evaluation Challenges, or
Create New Ones
1i. What can log analysis reveal
beyond outcomes?
Agent success or failure cannot be fully understood by outcome metrics
alone. Trajectory analysis can reveal validity challenges (Section 1a),
benchmark quality issues (Section 1b), contamination (Section 1f), and
surface underlying patterns in agent performance.
1j. What is the value of standardized
evaluation infrastructure for agents?
Model benchmarks were far simpler and cheaper. Trying to conduct
rigorous agent evaluations requires significant engineering effort and an
ecosystem of infrastructure such as reproducible environments, reusable
scaffolds, and common reporting standards. Standardized evaluation
infrastructure can help to simplify costs and develop greater
comparability of work.
1k. How should red-teaming evolve
for agents?
Model red-teaming was primarily focused on scrutinizing harmful
outputs and possible prompt vulnerabilities. For agents, there is a much
larger attack surface in addition to the model itself (e.g., from the
scaffold). Moreover, harmful model outputs in the agentic paradigm can
more directly lead to harmful action thanks to tool use and expanded
capabilities.
1l. What can open-world evaluations
reveal that benchmarks cannot?
Benchmarks flatten performance through scoped tasks and outcome
taxonomies. This was sufficient for the uni-dimensional form of model
outputs. But as expectations for agents rise in real-world settings,
understanding their performance demands more open-ended
examination, through detailed real-world cases with long-horizon,
complex agent trajectories. What does comparability look like for
open-world evaluations?
1m. How can evaluators interpret
agent performance relative to human
capabilities?
For model evaluations, comparison to human capabilities largely revolved
around comparing accuracy. But we are now demanding much more from
agents, applying them to work-like tasks. This creates new avenues for
comparison, by quantifying agent performance in terms of human time,
quality of output, and reliability.
Expanded Research Agenda
2a–2d. Human-agent interaction:
evaluation across simulated,
laboratory, field, and observational
settings
Model evaluations largely made sense in the context of prompts and
outputs. Agents are often situated in particular human-agent interactions
that are multi-turn and stateful, and require study across a variety of
evaluation settings: through simulated users, in laboratory settings with
humans, in field experiments, and through observational data on
real-world usage.
2e. Domain-specific evaluation
Model evaluation in specific domains largely tested whether models
possess specific knowledge within a domain. Agents are being tested on
domain-specific workflows.
2f. Economic impact evaluation
Model capability scores alone do not establish economic impact; actual
adoption, usage, and aggregate outcomes with agents are what matter.
2g. Agent behavior under evaluation
Agents in stateful environments create expanded opportunities to assess
the implications of reward hacking behavior and consequences of
evaluation awareness.
2h. Dangerous capabilities and misuse
Agent access to tool use, permissions, and stateful environments
amplifies ability to conduct consequential harmful actions.
2i. Agent and evaluation security
Security is no longer about the model alone, but now includes the
scaffold and wider evaluation environment as well.
5
Methodological Question or Area of
Evaluation
How Agents Exacerbate Existing Model Evaluation Challenges, or
Create New Ones
2j. Interaction harms
Model evaluation surfaced potential harmful outputs, but largely in
atomic individual responses. Persistent interactions with stateful agents
equipped with memory can create feedback loops that can accumulate
harms and lead to new ones across long trajectories.
2k. Systemic evaluation
Model evaluation was largely done in isolation. As agents embed in
organizations and institutions and interact with humans and other agents,
their effects on incentive structures demand study.
Scope and methods
Intended audience
This paper is written primarily for researchers who are interested in understanding the current state of
research of AI agent evaluation or to enter the field. It may be particularly useful for researchers who are
familiar with model-level benchmarking, but who are less familiar with agent-specific research. Researchers
who are already working on AI agents may find this paper helpful for understanding other aspects of the
broader agent ecosystem, or as an opportunity to reflect on the future direction of the field. We expect our
paper to be most relevant to researchers seeking generalizable knowledge, though much of the discussion
will also be useful to engineers doing downstream evaluations.
The paper also targets a broader audience of stakeholders interested in weighing the claims around AI
agents and their capabilities. For these readers, it offers a reference point for judging the current state
of agent capabilities, broadly and within particular domains. It does so by surfacing shared evaluation
challenges and emerging best practices within the latest paradigm of agents.
Table 2 suggests where readers with different goals might start.
Agents as the object of evaluation
We focus on evaluations of agent systems (“AI agents”): systems combining large language models (LLMs)
with scaffolds and tools to plan and act to achieve goals, without human intervention required at all
steps [129, 133, 142, 146, 246, 287, 305]. Moreover, we primarily discuss frontier AI agents whose
performance approaches the state-of-the-art across a wide range of capabilities [43]. We examine agent
evaluation at three levels of analysis: (i) the agent system itself, consisting of the combined behavior of its
components; (ii) agents in tandem with each other and with humans, e.g., human-agent interaction and
collaboration; and (iii) agents and their effects on institutions and organizations.
This survey emphasizes trends relevant to agent evaluations and how they have evolved relative to model
evaluations. We largely discuss agent systems rather than the model internals alone. For example, we
include work on agent security that emphasizes scaffold, environment, and harness design (e.g., indirect
prompt injection). However, we generally exclude from scope work solely concerning model internals, such
as mechanistic interpretability research into Sparse Autoencoders (SAE). We also exclude from scope work
that primarily focuses on simulation fidelity using “digital twins,” [217] rather than the capabilities and
risks related to the agents themselves.
We use evaluation to describe works that “measure, characterize, or analyze properties of AI models or
systems.” [43] Evaluations may occur before or after deployment, and may be conducted by a wide range
6
Table 2: Relevant sections organized by use case.
Use Case
Start With
Also Relevant
Entering
agent-evaluation
research
Agents as the object of evaluation; Validity (1a);
Part 2 introduction; Who evaluates (3a);
Disciplinary entry points (3i)
Venues (3g); Part 2 research
directions
Implementing an
evaluation
Metrics beyond accuracy (1c); Statistical
rigor (1d); Causal attribution (1e); LLM judges
and monitors (1h); Log analysis (1i); Evaluation
infrastructure (1j); Security (2i)
Contamination (1f); Behavior under
evaluation (2g); Rising costs (3f)
Weighing evaluation
claims
Validity (1a); Benchmark quality (1b); Statistical
rigor (1d); Contamination (1f); Human
baselines (1m)
Behavior under evaluation (2g);
Security (2i); Independence (3b);
Labs shaping evaluations (3c)
Evaluating agents for a
deployment
Field experiments (2c); Observational usage
data (2d); Domain-specific evaluation (2e)
Open-world evaluations (1l);
Simulated users and lab
experiments (2a–2b); Interaction
harms (2j)
Policy and oversight
Dangerous capabilities (2h); Systemic
evaluation (2k); Independence (3b); Public
sector (3h)
Red-teaming (1k); Economic
impact (2f); Labs and evaluations’
relationship (3c–3d)
Funding evaluations
Part 2 introduction; Independence (3b);
Funding (3e); Rising costs (3f)
Evaluation infrastructure (1j); Labs
and evaluations’
relationship (3c–3d); Venues (3g)
Model evaluation
Prompt
Model
Response
Grader
Reference answer
Agent evaluation
Agent system
Model
Scaffold
Tools
Task
Environment
actions
observations
Human
may intervene
at some steps
Trajectory
Final state
every step
Grader
Tests
Rubric
LLM judge
Figure 2: Model evaluation grades a single response against a reference answer. Agent evaluation covers a system of
LLM, scaffold, and tools that acts in an environment over many steps, sometimes with human intervention; instead of
a reference answer, the grader uses tests, a rubric, or an LLM judge to assess the trajectory (the process), the final
state (the result), or both.
7
of actors, including agent developers, governments, non-profit, or for-profit actors. Both technical and
sociotechnical forms of evaluation are in scope, with one complementing the other—technical evaluation
provides insight into capabilities and constraints of these systems, while sociotechnical evaluation allows us
to contextualize agents and understand how they interact in deployment settings.
We draw a few conceptual boundaries between evaluation and overlapping concepts of verification and
auditing. Evaluation is broader than verification work [43], since evaluators have methodological freedom
in designing and conducting evaluations, rather than strictly reproducing an existing process [255]. We
consider evaluation to be broader in purpose than auditing, since auditing incorporates evaluations primarily
in service of accountability [8, 39, 57, 198, 224]. We do not restrict our focus to evaluations with an
explicit accountability focus, but instead include it as part of the broader category of evaluations.
Methods
We use the term trend to describe a documented change in how agents are evaluated or in what evaluations
try to measure. Trends take on different forms—they may be new formulations of old problems or challenges,
for example in examining how agents create new validity challenges, or in new practices that have grown
to meet the needs within agent evaluations, for example in the increased reliance on LLM-as-judge setups.
We draw upon different forms of evidence, from public leaderboards, benchmarks, preprints, research
notes, online essays, blog posts, and position papers, from a variety of sources including industry, academic,
governmental, for-profit, non-profit, and philanthropic actors. Our examination of the ecosystem relies upon
snowball sampling, following citation chains, and examining points of collaboration between authors and
organizations. We rely solely on public-facing materials and publicly disclosed funding and grantmaking,
which at times limits the granularity of detail we can provide.
We structure our overall findings in three parts: the how, what, and who, of evals. While several trends
straddle multiple categories, our goal is to locate them alongside relevant trends expanding upon similar
themes. Our selection of trends is opinionated and the selection of citations within each trend is illustrative—
neither is intended to be comprehensive.
Part 1: Improved Evaluation Science and Methods (“How”)
Amid growing calls for a science of evaluations [26, 112, 113, 218, 293], AI agents have revealed important
shortcomings in previous evaluation practices, creating increased demand for scientific rigor [199]—in
how measurement is defined, how evaluation is conducted, and how it is reported. At the same time,
researchers increasingly recognize AI evaluation as a social science measurement challenge [282], one that
requires being explicit about constructs, how they are operationalized, and how they are measured.
This section begins by highlighting some animating concerns within the science of evaluations, by exploring
questions around validity, statistical rigor, causality, contamination and saturation. Next, we explore the
emerging toolkit for evaluations which aim to address these concerns: dynamic benchmarks, log analysis,
and LLM-based judging and monitoring (see Figure 3). Lastly, we examine the new evaluation formats
that are being applied to understand agents, as evaluators seek to complement benchmarks with new
approaches.
1a. How does validity change for agents?
Consider the example of a “replay agent” for computer use—one that merely copies a single successful
agent trajectory on a task and reproduces those exact actions [59]. No matter the state of the environment,
8
Construct
validity
Quality and
saturation
Statistical
rigor
Causal
attribution
Data
contamination
Ecological
validity
Design practices
1b
Item-level analysis
1d
Repeated runs
1e
Factorial designs
1f
Held-out test sets
Tools and formats
1g
Dynamic benchmarks
1h
LLM judges and monitors
1i
Log analysis
1j
Standardized infrastructure
1l
Open-world evaluations
addresses the concern
introduces a risk for it
Figure 3: How the methods discussed in Part 1 relate to its recurring concerns: which concerns each method is
presented as addressing, and where it is flagged as introducing a new risk.
the replay agent plods along. There is no reasoning involved, only blind recitation of a particular action
sequence. It scores highly on certain benchmarks, even achieving higher success rates than agents powered
by frontier models. Does this higher success rate imply that the replay agent is better at computer-use
reasoning than those of frontier models? What does “better” mean in this context?
This example highlights the importance of construct validity, which asks to what extent a metric—in this
case, success on various computer-use tasks—accurately represents the abstract concept (the “construct”)
it claims to measure: computer-use capability. This question of validity also assumes that the construct is
adequately defined, a prerequisite which itself may not be satisfied. Moreover, for agents, constructs of
interest are harder to operationalize due to the dynamic nature of their real-world interactions. Assessing
construct validity therefore requires a clear definition of both the construct of interest and its corresponding
metric [234].
The increased complexity of agents, along with their environments, scaffolds, and trajectories, has made
the task of assessing not only construct validity, but validity generally, more difficult and urgent. Construct
validity is one of many interrelated validity concerns (see Figure 4)—we discuss many of these in further
detail throughout this section, as validity underpins many trends as a first-order concern in evaluation.1
Internal validity specifies the strength of the causal relationship between changes to agents and measured
outcomes, and to what extent that attribution is defensible—a challenge we explore in Section 1e. External
validity examines whether evaluation results are generalizable to other contexts. We explore a specific
iteration of this concern, ecological validity in Section 1l, through the need for naturalistic evaluation
environments and settings.
best practices
• Define explicitly the underlying construct before choosing a metric. Walk through different ways that metric
1See also Goodhart’s law, which describes how a metric can diverge from the actual construct of interest.
9
Change to the agent
model, scaffold, tools
Metric
task success rate
Construct
computer-use capability
Other contexts
tasks, settings, users
internal validity
did the change
cause the difference?
construct validity
does the metric
capture the concept?
external validity
do the results generalize?
Ecological validity: do they hold under
real-world conditions of use?
Figure 4: Kinds of validity and the link each one concerns, illustrated with the computer-use example.
might be gamed by an agent strategy which lacks the capability of interest.
• Scrutinize the validity of evaluation results and the claims they support (e.g., using a validity frame-
work [234]), and how they rely on construct, internal, and external validity.
• Apply appropriate tooling (e.g., log analysis) to inspect possible sources of validity error—for example,
through shortcuts, contamination, or reward hacking [186]. We explore these methods in greater detail in
Sections 1h and 1i.
1b. How should we evaluate benchmark quality?
Underneath the broad umbrella of validity, a key concern is around the quality of an agent benchmark
as a whole. Ensuring benchmark quality is difficult for a few reasons. On a purely technical basis, tasks
are tedious for humans to check manually. Agents have amplified this challenge further, by expanding the
number of components where possible errors may reside. Moreover, constructs of interest are harder to
operationalize because of the dynamic nature of agents’ real-world interactions—lengthening the distance
between upstream evaluation and the downstream deployments of these systems. And with increasing
focus on multi-agent systems, the potential variance quickly becomes much larger.
Compounding this technical challenge are disincentives for benchmark quality assurance. Benchmarks
have touted how “impossible” their tasks are [231], while simultaneously overclaiming on the construct
a benchmark intends to represent [181]. Evaluation as a whole is also increasingly costly [92]: since
additional time and resources spent verifying tasks carry additional cost, benchmark developers may feel
pressure to prioritize other aspects of benchmark design. Benchmarks are also not immune to the pressure
to publish first—among similar benchmarks, the first one to be released has a strong advantage in adoption
by the community. While the peer review process might ideally surface potential issues, reviewers may
lack the time, resources, and expertise to conduct quality control on specific task contents, as specific task
contents are not within their remit.
Indeed, many potential problems are difficult to discern and become most visible once a benchmark reaches
saturation [186]. Saturation refers to when a benchmark can no longer reliably distinguish between the
performance of different agents (see Figure 5). For example, this may occur when several agents have
approached the benchmark ceiling and their performances have plateaued so closely together that they
are statistically indistinguishable from each other [6]. Saturation can be evidence of genuine progress in
underlying capabilities, while also offering an opportunity for introspection and iteration on a benchmark
and to identify validity issues. Following saturation, there is an opportunity to examine benchmark validity.
For example, plateauing performance could be indicative of reward hacking, or that the agent has exploited
some aspect of the evaluation test set [186, 268]. In the case of the replay agent, its successful performance
10
demonstrated the validity issues with its benchmark, a claim that would not necessarily merit attention
had its performance not reached parity with other approaches.
Successive agent releases
Benchmark score
benchmark ceiling
statistically
indistinguishable
Responses to saturation
Scrutinize the benchmark’s validity
Build a harder successor benchmark
Refresh the task set with new content
Keep the tasks that still discriminate
Incorporate metrics beyond accuracy
Figure 5: Schematic of benchmark saturation. Once the leading agents are statistically indistinguishable near the
ceiling, the benchmark no longer discriminates them, which prompts the responses on the right.
One response is to focus attention on developing the next, typically more challenging, version of the
benchmark. In fact, this very approach to iterating upon benchmarks can help set concrete targets
for future agent development (e.g., SWE-bench [124] and SWE-agent [302]). Another response is by
refreshing the task sets with new content, perhaps drawing from production traffic on a continual basis
(see Section 1g) [301]. Individual benchmark tasks may also be reused for greater benchmark efficiency,
through item-analysis of existing tasks to identify which tasks still discern between agents. For example,
a psychometric framework such as Item Response Theory (IRT) could be used to streamline the tasks,
such that a reduced subset of tasks still discriminates among top agents [191]. Alternatively, the same
benchmarks can be repurposed with a focus on new metrics beyond accuracy.
best practices
• Use saturation on one metric as a trigger for deeper benchmark scrutiny.
• Check for the possibility of impossible tasks, incorrect grading, task errors, and usage of shortcuts to review
benchmark validity.
• Apply item-level analysis (e.g., Item Response Theory (IRT) from measurement science) to determine which
benchmark tasks continue to discriminate among top agents, possibly to repurpose subsets into more efficient
benchmarks [191].
1c. What should agent evaluations measure beyond accuracy?
Even if task success serves as a valid metric, it may only capture one dimension of agent performance.
Evaluation has historically defined success in the context of a narrowly bounded task, for example through
a metric of “functional correctness.” [309] As agents take on longer trajectory tasks and deploy into
economically relevant areas of work, they merit a more holistic performance profile, adding, for example,
measures of cost, latency, or reliability.
This idea of bringing together many different measures of agent performance has roots in prior work on
model performance, such as the Holistic Evaluation of Language Models (HELM) [149]. Agents, which may
require multiple reasoning steps, access to tool use, and include much greater variance in cost and compute,
create additional dimensions for consideration. Relevant examples of quantitative metrics include:
11
• Cost: how many credits or compute an agent consumes. Cost can be a decisive factor in domain-specific
implementations, perhaps with respect to profitability [132, 133].
• Latency: how long it takes for models (or agents) to return a response. For time-sensitive applications,
latency has a substantial impact on usability. With agents, assessing latency requires not only factoring
in the model latency but also the latency of its other components (e.g., tool use) [133, 279].
• Reliability: a composite metric capturing how consistent, robust, and predictable an agent is [223].
Agents shipped in differing deployment scenarios may require adaptive thresholds of acceptable
reliability [223].
These quantitative metrics can also be complemented by qualitative analyses of agent trajectories, a trend
we explore further in Section 1i. Even if final outcomes are identical across multiple runs, the possibility of
highly disparate trajectories, perhaps with unsafe behaviors, merits study.
The way in which these dimensions together are being communicated also frames discussion around
agents. For example, Pareto frontiers [132, 133, Sec. 2.2, 82], such as between accuracy versus cost, or
performance versus token usage [31], more clearly enable joint consideration of these metrics [72, 161,
206, 286]. They help evaluators visualize the tradeoffs involved, since maximizing on certain metrics
such as accuracy alone can lead to unbounded cost [133]. This is especially relevant for agents, since cost
variance can be much higher for agents than models alone. These metrics may also be reported with respect
to release date to contextualize progress over time [114, 145, 223, 240].
Cost
Accuracy
on the Pareto frontier
dominated
Far along the frontier, large increases in
cost buy little additional accuracy.
Dominated: another agent is both
cheaper and more accurate.
Figure 6: Schematic accuracy vs cost plot. Agents on the Pareto frontier are not beaten on both dimensions by any
other agent; the remaining agents are dominated.
best practices
• Report accuracy jointly with an appropriate notion of cost (e.g., monetary value, latency) and illustrate the
tradeoff as a Pareto curve [133] (see Figure 6).
• Measure reliability as a metric in its own right and set an acceptable threshold scaled to the specific
deployment scenario.
• Clarify what model-scaffold combinations are being measured, to better understand model × scaffold
interaction effects.
1d. How can agent evaluations achieve statistical rigor?
Agent performance on benchmarks can vary widely, due to particular tasks selected into the benchmark, or
variance across different agent trajectories within the same task. Certain approaches to aggregating or
12
subsetting these trajectories may obscure this variance. Increased statistical rigor can capture this variance
and help to address concerns around cherry-picking in measuring agent performance, and distinguish
between best-case and typical performance, which may be of interest to different audiences and interpreted
differently as such.
Repeated runs are increasingly adopted for agent evaluations, but single-run evaluations remain common due
to the constraints from cost and compute [16, 223, 255]. One popular performance metric is pass@k [49],
originally popularized for measuring functional correctness of AI-generated code [49] and subsequently
adopted for agent evaluations [132, 133, 278, 302]. Pass@k represents the probability of “success” at least
once within k attempts. For example, pass@1 communicates the expected performance of an agent on a
single attempt [40]. Yet, for stateful interactions such as computer use, pass@k may be an inappropriate
measure for reliability or consistency [59]. In contrast, pass∧k quantifies consistency in task success by
representing the probability of success over all k runs [304] (see Figure 7). Neither metric necessarily
captures how robust an agent is to changes in the testing environment; for example, gaming through a
replay agent (see Section 1a) may still report high scores for both metrics while still being very brittle for
other environments. Multi-agent contexts introduce additional sources of variance evaluators must contend
with.
Five runs per task
filled: success
hollow: failure
pass@1
a single run
succeeds
pass@5
at least one run
succeeds
pass∧5
all five runs
succeed
Task A
5/5
Task B
3/5
Task C
1/5
Task D
0/5
Across the four tasks
45\%
75\%
25\%
Figure 7: pass@1, pass@k, and pass∧k computed on the same illustrative runs (k = 5). pass@k credits a task after
a single success, whereas pass∧k requires every run to succeed.
Reporting point estimates alongside uncertainty with error bars is a standard practice among many
evaluators [182, 279], but evaluations are broadening to specify sources of variation as well. Generalizability
theory, rooted in the field of measurement science, has a long history of formalizing this need for variance
decomposition—and can serve as a framework for tackling these challenges in agent evaluation as well.
One important distinction for rigor in evaluation is between benchmark accuracy—which concerns expected
performance of an agent on a specific benchmark—and generalized accuracy—which concerns the tasks that
get selected into a benchmark relative to the entire population of tasks. That is, benchmark accuracy asks,
“how well would an agent perform on these exact tasks again?” In contrast, generalized accuracy asks, “how
well would an agent perform on new tasks similar to these tasks?” The latter has wider uncertainty because
sampling the tasks themselves adds variance. This distinction between benchmark and generalized accuracy
also has clear roots in measurement science, through latent variable modeling. NIST has explored two
approaches to estimating benchmark accuracy and generalized accuracy. First, a regression-free approach
directly averages across task-level success rates without quantifying capability or difficulty. Second, a
generalized linear mixed model (GLMM) approach estimates benchmark accuracy by averaging over tested
tasks, and estimates generalized accuracy by generating a distribution of possible task difficulties and
13
averaging over said distribution. As another example, the HiBayES framework for hierarchical Bayesian
models from the UK AI Security Institute (UK AISI) factors in the families of tasks and domains that can
characterize subdivisions of a larger benchmark [157, 273]. These approaches involve tradeoffs between
efficiency and relying on various assumptions about the benchmark, namely, how tasks are sampled for the
benchmark and the difficulty of those tasks [134].
best practices
• Conduct repeated runs to measure consistency for each task, and report how many runs were conducted.
Avoid claims relying on low run-to-run variance if only a single run is affordable.
• Ensure that the metric of interest and the stated claim are aligned: report pass@k when it matters whether
an agent can succeed at all, and report pass∧k for whether it succeeds every time—pass@k should not serve
as evidence of reliability.
• Account for different sources of variance (e.g., judge error) in reported uncertainty, particularly for LLM-as-
judge setups (see Section 1h), where raw scores may understate uncertainty.
1e. How can agent evaluations attribute causal changes in performance?
If an agent’s performance does improve in a statistically significant manner, can we pinpoint why? Agent
performance is jointly influenced by the model, scaffold, tool use, budget constraints, the environment,
and how it is evaluated [133]. When several of these components change at once, a statistically significant
impact on performance does not allow us to disaggregate which precise factors drove that change (see
Figure 8). And even when only one component changes, its effect may be attributable not only to the
change in isolation, but rather to interactions with the system as a whole.
For example, an agent scaffold may create variable effects when paired with different models. As the
PaperBench [252] benchmark highlights, some configurations may yield improvements on metrics, while
others may not. This further complicates this question of attribution, and highlights the need to think of
the interaction between model and scaffold in a holistic manner—that is, as a model × scaffold interaction
effect.
(a) Changing both at once
Model A
Model B
Model C
Scaffold 1
Scaffold 2
Scaffold 3
The two runs differ in model and scaffold, so the
difference cannot be attributed to either one.
(b) Factorial design
Model A
Model B
Model C
Scaffold 1
Scaffold 2
Scaffold 3
Every pairing is run with environment and budget
fixed, separating model effects, scaffold effects and
their interaction.
Figure 8: Attributing a change in agent performance to model vs scaffold. (a) When model and scaffold change
together, their effects are confounded. (b) A factorial design runs every model-scaffold pairing under fixed conditions.
Carefully parsing these factors helps to ensure that causal claims are well supported. Even though agents
are built around a core underlying LLM, improvements or limitations of an agent cannot be used to
infer limitations of a model itself. It is also worth asking whether existing limitations are attributable
14
to implementation challenges such as under-elicitation [140, 169].2 Indeed, in work such as CORE-
Bench [248], significant gains in performance were made possible by improved scaffolds, even when models
were held constant.
Disaggregating these different components enables greater comparability over time, as models, scaffolds, and
environments evolve independently. The growing need to test in multi-agent settings further complicates
the work to achieve comparability. Attribution therefore requires instituting appropriate controls and
baselines in the experimental design in order to best isolate specific effects; for example, through keeping
budgets and environments fixed, and systematically varying discrete components [132, 300]. And in
continuity with model benchmarks, systematically varying structured prompt templates for agents can help
to facilitate the robustness of causal findings [183].
best practices
• Identify and document clearly different components of an experimental setup—the models used, scaffolds
paired, environments tested, prompts used, and the budget constraints. When comparing agents, hold all
but one component fixed such that any change in performance can be properly attributed.
• Do not attribute an agent’s failure directly to the underlying (language) model, without accounting for
whether improved scaffolding around the model changes the performance profile [140, 248].
• When budget allows, systematically test model × scaffold interaction effects in various permutations, e.g.,
with (fractional) factorial designs—in which multiple independent variables are tested in conjunction.
1f. How can agent evaluations resist contamination?
Evaluating agents involves greater challenges of contamination relative to model evaluations alone. For
model-only evaluations, contamination arises from exposure to evaluation tasks verbatim or with high simi-
larity, through pre-training data, benchmark content, or public test sets [42, 119, 130]. Evaluation datasets
may also leak through direct and indirect channels (e.g., via synthetic data), and lead to overfitting [164].
Agent evaluations introduce additional surfaces for contamination (see Figure 9). Test sample knowledge
may end up directly programmed into the scaffold, prompts, tool use, or other components of the agent
system, even without the scaffold developer’s knowledge, a problem described as “harness-level cheat-
ing” [253]. With access to tool use, agents may be susceptible to search-time contamination, whereby agents
are able to retrieve answers at inference time. Moreover, longer agent trajectories create more opportunities
for potential leakage and raise the costs of inspection.
Training data
benchmark tasks seen
verbatim or with high
similarity during pre-training
Scaffold and prompts
knowledge of test samples
built into the scaffold,
prompts or tools
Tool use at test time
answers retrieved at inference
time (search-time
contamination)
also present in model evaluation
additional surfaces for agents
Figure 9: Stages at which contamination can enter an agent evaluation.
Contamination should also not be conceptualized as merely a binary question of whether certain data
is included or not. Between direct memorization and robust generalizability lies a spectrum of more
2Elicitation refers to “getting maximal performance of a model or system on a capability evaluation” [111, 169, 266].
15
subtle forms of lack-of-generalization [133, Sec. 5, 54, Sec. I.3.2], in which agents to various degrees may
exploit knowledge or structural conventions that do not map to environments the benchmark is intended
to measure. Consequently, evaluation should specify the intended level of generalization that a benchmark
aims to test and design accordingly.
best practices
• Define secret holdouts appropriate to the level of generality expected of the agent [132], and limit repeated
use and optimization against the holdouts [46, 104, 106].
• Report test-train overlap [307], any benchmark-specific information used to develop the agent system’s
prompts, scaffolds, and tool usage, and create transparency on private testing practices [249].
• Inspect agent trajectories (see Section 1i) for contamination and incorporate clean-room evaluation harnesses
into the testing setup [126].
• When a holdout set is publicly available, update the set regularly with refreshed tasks (Section 1g) such that
benchmark scores do not depend on tasks the agent may already have been exposed to.
1g. What role do dynamic benchmarks play, and what are their tradeoffs?
The term dynamic benchmark is currently an unsettled category and encompasses a broad set of different,
non-exclusive approaches. Here, we use dynamic benchmarks as an umbrella term to include evaluations
whose tasks, environments, or labels change over time or through changes in model or agent behavior.
Dynamic benchmarks are relevant for agents because they rely not only on an underlying model, but also
on scaffolds, tools, and environments, all of which are liable to change over time, and which the agents
themselves may change.
For example, dynamic agent benchmarks may:
• Provide extensible environments that allow new tasks to be updated or modified (e.g., AgentDojo) [62];
• Evaluate agents on “live” data with test sets that are regularly refreshed [120, 299], perhaps using
production traffic [313];
• Compare predictions against new information to test forecasting ability [306];
• Focus on adversarial data collection, where new tests are selected from previous model failures, e.g.,
Dynabench [70, 136];
• Evaluate agents through competition against adaptive agents, e.g., Agent Island [185].
Each of these approaches offers particular responses to certain limitations of static benchmarks. Refreshing
or updating benchmarks with new tasks can aim to mitigate possible contamination—for example, LiveBench
adds new test questions on a monthly basis to resist contamination [299]. Drawing from production traffic
can aim to increase ecological validity. Adversarial data collection can respond to saturation by probing for
new failure modes.
A central tradeoff with modifying and refreshing tasks over time raises the issue of longitudinal compara-
bility; if the evaluation itself is not consistent across different tests over time, how can we compare the
results in a relative sense? While comparability is possible when agents are evaluated against the same
version of a dynamic benchmark, evaluations across different versions may not be comparable. Looking
toward the field of psychometrics may offer solutions to this issue by embracing test equating and linking,
whereby tests are designed such that scores on different versions are interchangeable through a shared
scale (see Figure 10).
Validity remains a first-order concern—just because a benchmark is dynamic does not necessarily ensure
16
Version 1
Version 2
Version 3
anchor tasks
refreshed tasks that differ in each version
Scores on a shared scale,
comparable across versions
Figure 10: Test equating for a dynamic benchmark. A common set of anchor tasks links versions whose other tasks
are refreshed. This allows us to place scores on a shared scale.
the test set continues to accurately map to the construct of interest. For example, in forecasting applications,
the set of sampled events needs adequate justification for whether it represents the intended forecasting
domain. Adversarial data collection approaches such as Dynabench intentionally concentrate the data
toward failures. Although this can help to surface specific failure modes that other benchmarks gloss over,
it can also mean that there is a distributional shift away from more realistic usage patterns. Consequently,
performance on such benchmarks may overweight edge cases relative to the conditions under which models
are typically deployed [87, 107, 247].
best practices
• Apply test equating and linking from psychometrics to ensure a common anchor subset across scores, for
longitudinal comparability.
• Track distribution shifts and regularly revisit whether results validly support claims about the intended
construct.
• Implement version control for each component of the benchmark, to enable reproducibility and comparability
over time.
1h. Where should LLMs judge and monitor agents?
LLM-as-judges (or “autograders”) are becoming a core component for evaluation for open-ended outputs
and in scenarios where human expert graders are costly [103]. For example, OpenAI’s PaperBench develops
outcome-focused judging rubrics for its LLM-based judge, SimpleJudge, and then validates its judgement
of the final submissions against human expert labels using another evaluation, JudgeEval [252]. In recent
work, UK AISI has used a panel of LLM graders from which a median score is reported [221]. LLM-based
judges are also playing a key role in understanding not only outcomes but agent trajectories, for example
through the Agent-as-a-Judge framework [312].
However, this can come with challenges around grader reliability. In effect, who judges the judges? This is
a key question raised when LLM-as-judge setups become integrated into the experimental setup with its
potential to introduce “biases and adversarial vulnerabilities” into the verification process [59]. Relying on
these tools therefore creates another surface requiring evaluation. LLM-as-judges require validation against
appropriate holdouts, such as expert labels. Researchers have also examined evaluating LLM-as-judge
setups using crowdsourced human preferences [308], or through evaluation against response pairs [261].
Another response is to model judge error rather than agreement alone—for example, UK AISI has also
systematically assessed the quality and quantified biases in its LLM-as-judges using generalized linear
models (GLMs) [68]. This judge error should moreover be reflected in how uncertainty is reported, in
particular to inform how error bars are presented when LLM-as-judges are in play.
Alongside LLM-as-judge setups are LLM-based monitors [32] (or “scanners”) which inspect agent trajectories,
17
either in real time or after the fact, to flag specific instances of agent behavior. In particular, current
approaches focus on monitoring the agent chain of thought (CoT)—the visible, intermediate reasoning
steps returned by the model within an agent. Monitoring can be valuable for detecting reward hacking,
cheating, or particular unsafe behaviors from agents. However, CoT is not foolproof; using this technique
may not fully surface sandbagging or other scheming behaviors—in effect, the absence of evidence may not
imply evidence of absence [239]. Moreover, LLM-based monitoring also carries risks of overreliance. The
volume of agent trajectories, particularly in multi-agent contexts, is outpacing the ability of evaluators to
manually review this content—a point emphasized in METR’s investigation into the OpenAI–Hugging Face
incident [179]. As manual analysis becomes less feasible, the need for validating monitoring approaches
and quantifying their respective biases becomes even more pressing. We explore these challenges and
different evaluation targets in Section 2g.
best practices
• Validate autograders and scanners against human-coded holdouts representative of the intended use case
with subject matter expert participation.
• Incorporate targeted manual analysis into automated analysis. In particular, review a subset of the cases a
judge passed/stopped to quantify how many cases might slip by undetected.
• Quantify biases in autograder setups, through human-LLM / human-agent agreement and monitor if/when
training traffic is no longer representative of the object of judgement.
1i. What can log analysis reveal beyond outcomes?
Outcomes alone are insufficient for examining agent performance, as they compress an agent’s trajectory
into metrics that may not fully explain why an agent was successful or not (see Figure 11). Logs—the
records generated when using an AI system that may include model inputs and outputs, environmental
interactions, and metadata—offer an opportunity for analyzing evidence from agent trajectories [69, 272].
Log analysis can prove useful for ensuring construct validity and other potential measurement validity
concerns. For example, LLM-aided log analysis for the Holistic Agent Leaderboard (HAL) surfaced instances
of reward hacking and contamination as an agent searched for benchmark answers on Hugging Face [132].
Log analysis can also assist in addressing challenges to (1) internal validity, (2) external validity, and (3)
safety evaluation [140, Table 1]. Log analysis enables identification of failure modes, instances where
shortcuts have been taken, or reward hacking [133, Sec. 5].
Both manual review and automated analysis are complementary to each other. Manually exploring a log
sample can help evaluators develop a rubric for later automatic analysis or monitoring. Manual review
can also help to reveal quality issues hidden by aggregate scores. For example, in METR’s review of its
evaluation methods, manual review was able to identify that algorithmic scoring had overinflated the
performance on certain tasks and obscured reward hacking by accessing the answers via GitHub [173].
However, log analysis has its limitations. First is a question of scale. The sheer volume of agent transcripts,
particularly in the context of multi-agent interactions, limits the feasibility of manual review and i
\end{document}
